\subsection{Initializing one more puncture by products of fibre primes}

\begin{statement}{A.83}{Proved}
{An all-place initialization with prescribed local points}
Let $4\mid n$, let $c_1,\ldots,c_k$ be distinct rational numbers,
and let $a_i\in\Q^\times$.  Suppose the affine old cover
\[
 C:\quad T-c_i=a_iX_i^n\quad(1\le i\le k),
                         \qquad\prod_iX_i\ne0
 \tag{A.83.1}
\]
has a point over every completion of $\Q$.
Let $c\in\Q\setminus\{c_1,\ldots,c_k\}$ and put
\[
 D_i=\frac{c-c_i}{a_i},\qquad
 K=\Q(\sqrt{D_1},\ldots,\sqrt{D_k}).
 \tag{A.83.2}
\]
Assume $[-1]$ is outside the squareclass span of the $D_i$.
For a finite set of places $S_0$, prescribe points
$P_v\in C(\Q_v)$ with $T(P_v)\ne c$.
There is $a\in\Q^\times$ for which the augmented cover
\[
 C':\quad T-c_i=a_iX_i^n\ (1\le i\le k),
       \qquad T-c=aY^n,\qquad Y\prod_iX_i\ne0
 \tag{A.83.3}
\]
has a point over every completion of $\Q$, and each prescribed
point $P_v$ lifts to an actual point of $C'(\Q_v)$.
The forgetful map $C'\to C$ is the displayed map forgetting $Y$;
in particular the old global powerclasses are unchanged.
The coefficient $a$ can also be chosen to have valuation one
at a fresh prime outside any specified finite set, at which
all old coefficients and puncture differences are units.
\end{statement}
\begin{proof}
The proof uses products of primes with a fixed point on the fibre
at $c$.  We specify the ray-class condition and prove its
compatibility from the actual local evaluations.

\emph{The finite set and the local target classes.}
The geometric degree of the augmented cover is $N=n^{k+1}$:
the valuations of $T-c_1,\ldots,T-c_k,T-c$ at their distinct
zeros give independent classes modulo $n$th powers.
Its smooth projective completion has inertia of order $n$
above these $k+1$ finite punctures and above infinity.
Thus its genus $g$ and its number $b$ of geometric boundary
points satisfy
\[
 2g-2=-2N+(k+2)N(1-1/n),\qquad b=(k+2)n^k.
 \tag{A.83.4}
\]
These are the actual Kummer and ramification maps computed in
(A.48.4)--(A.48.6).  In particular the bound is independent
of the nonzero coefficients.  Fix an integer
\[
                         B>(2g+b+2)^2.
 \tag{A.83.5}
\]
Let $S$ be a finite set of rational primes containing the finite
places in $S_0$, every prime at most $B$, all prime divisors of
$n$, all primes in the numerators and denominators of the $a_i$
and of the nonzero pairwise puncture differences, all primes
in the denominators of $c_1,\ldots,c_k,c$, and every prime in
any additionally specified finite set.
Enlarge $S$ to contain the ramified primes of the splitting field
of $Z_i^n=D_i$ and the conductors of all characters of $K$.
All these are finite sets.

At each $p\in S$ choose an old point $P_p$ with $T(P_p)\ne c$,
keeping any prescribed $P_p$.  Choose $P_\infty$ in the same
way, retaining it when prescribed.  These choices exist:
the old cover has a point at each place; its equations have
nonzero derivatives in the $X_i$ at any affine point, and
their implicit-function parametrization by $T$ gives local
points with $T$ different from the one forbidden value.
Put
\[
 \alpha_p=T(P_p)-c,\quad \alpha_\infty=T(P_\infty)-c,
 \qquad
 d=\operatorname{sgn}(\alpha_\infty)
                          \prod_{p\in S}p^{v_p(\alpha_p)}.
 \tag{A.83.6}
\]
Then $d$ is a nonzero rational $S$-unit,
$u_p=\alpha_p/d\in\Z_p^\times$ for every $p\in S$,
and $\alpha_\infty/d>0$.

Choose integers $N_p>2v_p(n)$, increasing them so that
\[
 M=\prod_{p\in S}p^{N_p}
       \quad\text{satisfies}\quad K\subset R:=\Q(\mu_M).
 \tag{A.83.7}
\]
Quadratic conductors give the latter inclusion explicitly.
The Chinese remainder theorem gives a unit class
$\rho\in(\Z/M)^\times$ with
\[
                    \rho\equiv u_p\pmod {p^{N_p}}\quad(p\in S).
 \tag{A.83.8}
\]
Here the residues of $u_p$ are taken in $\Z_p/p^{N_p}$.
We also write $\rho$ for the arithmetic cyclotomic automorphism
$\zeta_M\mapsto\zeta_M^\rho$ of $R$.

\emph{Why this exact ray class fixes $K$.}
By A.80, every character of $\operatorname{Gal}(K/\Q)$ is
$\chi=\sum_i t_i\chi_{D_i}$, $t_i\in\{0,1\}$, and on the
old curve its cyclic algebra is the constant
\[
 (T-c,\chi)=\mathcal B_\chi,\qquad
                 \mathcal B_\chi=\sum_i t_i(a_i,D_i)
                                      \in\operatorname{Br}(\Q)[2].
 \tag{A.83.9}
\]
All finite ramified places of these constants are in $S$.
Using the local invariant--reciprocity comparison for cyclic
algebras, evaluation at the actual points $P_v$ gives
\[
 \sum_{p\in S}\chi(\operatorname{rec}_p\alpha_p)
       +\chi(\operatorname{rec}_\infty\alpha_\infty)
       =\sum_v\operatorname{inv}_v(\mathcal B_\chi)=0.
 \tag{A.83.10}
\]
The convention is arithmetic Frobenius on unramified uniformizers;
on cyclotomic roots the local reciprocity of a unit is its inverse,
as in \cite[I, Example 3.13; III, Proposition 3.6]{MilneCFT}.
Global reciprocity for the principal idele $d$ gives the same
zero sum with $\alpha_v$ replaced by $d$: outside $S$ the number
$d$ is a unit and $\chi$ is unramified, so those local terms vanish.
Subtract these two equalities.  The real term vanishes because
$\alpha_\infty/d>0$.  Hence
\[
                    \sum_{p\in S}\chi(\operatorname{rec}_p u_p)=0.
 \tag{A.83.11}
\]
In $\operatorname{Gal}(R/\Q)$ the product of these unit reciprocity
images is exactly $\rho^{-1}$: at $p$ it acts on the $p$-power
cyclotomic factor by $u_p^{-1}$ and on the prime-to-$p$ factor
trivially.  Formula (A.83.8) specifies all these factors.
Thus (A.83.11) says $\chi(\rho)=0$ for every character of $K$.
Since characters separate a finite abelian group, we have proved
\[
                            \rho\in\operatorname{Gal}(R/K).
 \tag{A.83.12}
\]
This deduction uses the prescribed local points themselves;
it does not assume an unproved constancy of an unramified class.

\emph{The exact stabilizer image and the prime product.}
Let $L$ be the splitting field of the finite fibre algebra
$\mathcal E=\Q[Z_i^{\pm1}]/(Z_i^n-D_i)$, and put
$G=\operatorname{Gal}(LR/\Q)$.
Let $\Omega$ be its root-tuple set and
$\pi_R:G\twoheadrightarrow\operatorname{Gal}(R/\Q)$ the
restriction map.  For every $\omega\in\Omega$ the fixed field
of $G_\omega$ is its actual root field
$E_\omega=\Q(z_1,\ldots,z_k)$.  Therefore
\[
 R^{\pi_R(G_\omega)}=R\cap E_\omega,\qquad
 R^{\langle\pi_R(G_\omega):\omega\in\Omega\rangle}
      =R\cap\bigcap_{\omega\in\Omega}E_\omega=K.
 \tag{A.83.13}
\]
For the last equality, the left-hand field is abelian over $\Q$
and is contained in every root-field factor; A.80 identifies
the largest such field as $K$, and $K\subset R$.
Consequently the images of the actual stabilizers generate
exactly $\operatorname{Gal}(R/K)$.

By (A.83.12), write
$\rho=\gamma_1\cdots\gamma_t$ with each
$\gamma_j=\pi_R(g_j)$ for some $g_j\in G_{\omega_j}$.
Inverses cause no problem because each individual stabilizer
image is already a subgroup.  Chebotarev, in the precise form
\cite[Theorem 8.31]{AgentMilneANT}, supplies distinct rational
primes $q_j$ outside $S$, larger than $B$, with Frobenius
conjugate to $g_j$.  Their restrictions to the abelian field
$R$ are exactly $\gamma_j$, so
\[
                   \prod_{j=1}^t q_j\equiv\rho\pmod M.
 \tag{A.83.14}
\]
Frobenius at $q_j$ fixes a root tuple: conjugacy preserves this
property.  Since the integral fibre is etale at these primes,
it therefore has a point over $\mathbf F_{q_j}$.
Choose also a distinct prime $q_*$ outside all the specified
finite sets that splits completely in $LR$.
It has the same fibre-point property and $q_*\equiv1\pmod M$.
Set
\[
                         a=d\,q_*\prod_{j=1}^t q_j.
 \tag{A.83.15}
\]
The extra factor ensures $v_{q_*}(a)=1$ at a fresh prime;
it does not alter the prescribed ray class.

The product in this step is essential even when the common
character field is trivial.  For $n=4$ and $D_1=25$, the
fibre polynomial factors as
$(Z^2-5)(Z^2+5)$, with root fields
$\Q(\sqrt5)$ and $\Q(\sqrt{-5})$ and common field $\Q$.
Modulo $20$, the ray class $13$ fixes no root: both quadratic
characters take the value $-1$ there.  The classes $3$ and
$11$ each fix a root and their product is $13$ modulo $20$.
Chebotarev realizes their classes by arbitrarily large distinct
primes.  Thus the generated stabilizer image, rather than its
union alone, is the exact group needed in (A.83.13).

\emph{All local points and the precise forgetful map.}
For $p\in S$, (A.83.8), (A.83.14), and (A.83.15) give
$a/\alpha_p\in1+p^{N_p}\Z_p$.
Every element of this group is an $n$th power: Hensel applied
to $Z^n-a/\alpha_p$ at $Z=1$ uses
$N_p>2v_p(n)$.  Thus $\alpha_p/a$ is an $n$th power as well.
Keep the exact old coordinates of $P_p$ and take any such
root for $Y$.  This lifts $P_p$ through (A.83.3).
At the real place, $a$ and $\alpha_\infty$ have the same sign,
so their ratio is positive and has an $n$th root.  The same
construction lifts $P_\infty$.

At a fresh coefficient prime $q\in\{q_*,q_1,\ldots,q_t\}$,
choose a nonzero residue root tuple
$\bar X_i^n=D_i$ supplied by its Frobenius fixed point.
Take
\[
                              T=c+a,\qquad Y=1.
 \tag{A.83.16}
\]
Here $v_q(a)=1$ and all $D_i$, $a_i$, and $c-c_i$ are units.
The equations
$X_i^n=(c-c_i+a)/a_i$ reduce to the chosen residue equations.
Their derivatives $n\bar X_i^{n-1}$ are units, so Hensel lifts
each root independently.  This gives an actual local point
with all coordinates nonzero.

At any other prime $p\notin S$, all coefficients, differences,
and $n$ are units, and $p>B$.
The reduction of (A.83.3) is the smooth geometrically connected
Kummer cover used in (A.83.4).  The Hasse--Weil estimate gives
at least
\[
                        p+1-2g\sqrt p-b>0
 \tag{A.83.17}
\]
interior residue points, by (A.83.5).
The trace and eigenvalue statements giving this bound are
\cite[(1.5.1), Theorem (1.6)]{AgentDeligneWeilI} and
\cite[Proposition 14.2]{AgentMilneLEC}, as in A.48.
An interior point has all coordinate derivatives nonzero;
Hensel therefore lifts it to a $\Q_p$ point of $C'$.
These cases exhaust the places of $\Q$ and prove every assertion.
\end{proof}

The result initializes one additional finite-exponent coordinate
in the specified squareclass range, with the exact old coefficients
preserved.  A.85 gives the complete finite-character criterion
for the complementary $[-1]$-span range as well.
Neither statement identifies the compatible limits of
successive initializations with an ordinary completed tuple for
all rational punctures.  The finite-level obstructions A.69 and
A.73 and the completed-tuple restriction A.76 retain their stated
scopes.

The finite certificate \path{code/puncture_initialization_certificate.py}
checks the displayed proper-union example by the literal fourth-power
equation at finite primes, its ray-class products, the genus and
boundary counts used in the uniform inequality, and the elementary
local power congruences.  It does not claim to execute the
arbitrarily large Chebotarev construction.
